% Evaluation section, WWW 2027 submission cut. ACM sigconf, anonymous review.
% Build: python3 eval/paper/build_eval_pdf.py  (tectonic fetches acmart on demand)
% Figures: eval/figures/www2027/, produced by eval/plot_professor_figures.py from
% eval/results/lan4d/summary_lan4d.json and eval/results/fabric_cloud.json.
\documentclass[sigconf,anonymous,review]{acmart}
\settopmatter{printacmref=false}
\setcopyright{none}
\renewcommand\footnotetextcopyrightpermission[1]{}
\acmConference[WWW '27]{Proceedings of the ACM Web Conference 2027}{April 2027}{Dublin, Ireland}
\usepackage{booktabs}
\graphicspath{{../figures/www2027/}}

\title{Knowledge-Aware Blockchain Transaction Validation with SHACL}
\subtitle{Evaluation section, WWW 2027 draft}
\author{Anonymous Author(s)}

\begin{document}
\maketitle

\setcounter{section}{4}
\section{Evaluation}
\label{sec:eval}

\subsection{Questions, treatments, and setup}
\label{sec:setup}

We ask four questions. \textbf{RQ1}: do the treatments reproduce the
intended marketplace rules, and does a reused verdict ever disagree with a
fresh one? \textbf{RQ2}: what does one validation cost, what does a cache hit
still cost, and does that hold as the ledger grows? \textbf{RQ3}: how do
lifecycle latency and achieved throughput differ end to end? \textbf{RQ4}:
which of these behaviors survive on Fabric, where SHACL runs only as a
pre-endorsement check? Comparisons stay inside one platform.

\paragraph{Treatments}
\emph{Direct} is the legacy imperative validator. \emph{SHACL} sends every
validation to a sidecar that evaluates the shapes and then runs the
ledger-state checks (ownership, effective advertisement status, linkage,
uniqueness) against MongoDB. \emph{SHACL+$\Delta$} adds a verdict cache keyed
by the transaction id and a hash of a declared \emph{dependency projection}:
for each transaction type, the dependency operation, asset and metadata
records, the effective advertisement status, and the identifiers returned by
the uniqueness queries. The projection is written per transaction type, not
extracted from the shapes, and it hashes whole records, so a change to an
unused field of a referenced record still misses. Every lookup recomputes the
hash, so safety never depends on a commit notification; notifications only
retire the committing transaction's own entry and drop entries indexed under
the entities it affects. The sidecar reports the transaction ids named in its
queries. During this campaign the validator counted reads outside the
projection and observed none; the current code additionally refuses to cache
such a verdict. That report names ids, not predicates, so a uniqueness query
that finds nothing is indistinguishable from one that was never declared.
Coverage is therefore a guard on these runs, not a proof.

\paragraph{Where validation runs}
On SmartChainDB a transaction is validated at HTTP admission in the web
process, then in CheckTx and DeliverTx in the ABCI process. The two processes
keep separate caches, and only ABCI sees commits. On Fabric the chaincode
decides; SHACL runs in the client before endorsement.

\paragraph{Setup}
Four validators with equal voting power, each 4~vCPU and 8\,GB, share a
private network in one region. Each runs Tendermint~0.31.5, MongoDB~3.6,
BigchainDB with one gunicorn worker, and the Node sidecar; a separate
2-vCPU client submits. Fabric~2.5 runs the \texttt{fabric-samples} test
network (two organizations, two peers, one orderer) with Go chaincode on one
host. The LAN campaign uses 30 seven-step lifecycles per treatment, 500
offers on one advertisement, 500 offers split around a close, 200 disjoint
offers, and 24 conformance cases repeated three times; in-process timings use
$n{=}200$ per condition. Intervals are percentile-bootstrap 95\% confidence
intervals of the median (2{,}000 resamples), written median [lo, hi]. Samples
within the campaign are not independent deployments. Fabric lifecycles use
$n{=}5$ and report min--max.

\subsection{Correctness and safe reuse (RQ1)}
\label{sec:safety}

\paragraph{Conformance}
The oracle is a truth table of the intended rules, fixed before the runs: one
valid case per operation and invalid cases for the wrong party, a missing
dependency, a bad status, a mismatched asset, and a duplicate return or
acceptance. SHACL and SHACL+$\Delta$ match all 72 executions with no server
error. Direct matches 66: it accepts a second \texttt{REQUEST\_RETURN} for a
sale and a second \texttt{SELLER\_ACCEPT\_RETURN} for a return, three times
each. Direct therefore does less checking than the SHACL treatments, and its
times below are a legacy reference, not a same-semantics baseline. All three
treatments commit the 250 offers placed before an advertisement closes and
reject the 250 placed after it; that 50\% abort rate is constructed. With
every sidecar stopped, both SHACL treatments reject 3 of 3 otherwise valid
offers and none commits (Table~\ref{tab:conformance}).

\begin{table}[t]
\caption{Conformance. The last two rows are the Fabric probes.}
\label{tab:conformance}
\small
\setlength{\tabcolsep}{4pt}
\begin{tabular}{@{}lccc@{}}
\toprule
Check & Direct & SHACL & SHACL+$\Delta$ \\
\midrule
Oracle matches (24 cases $\times$ 3) & 66/72 & 72/72 & 72/72 \\
Post-close offers rejected & 250/250 & 250/250 & 250/250 \\
Sidecar stopped: rejected, none committed & n/a & 3/3 & 3/3 \\
Fabric: non-owner advertisement rejected & yes & yes & yes \\
Fabric: post-close offers rejected & 10/10 & 10/10 & 10/10 \\
\bottomrule
\end{tabular}
\end{table}

\paragraph{Reuse under change}
Eight controlled sequences cover four dependency classes twice each. A
pending transaction is validated (cold miss) and validated again (hit); an
unrelated \texttt{CREATE} commits; then a commit changes its dependency
(closing the advertisement, a competing sale, a first return request, a
first acceptance). After each step we compare the served verdict with a fresh
engine verdict. Figure~\ref{fig:seq} shows three policies with and without
the commit signal. $\Delta$ keeps all 8 unrelated verdicts and, after the
relevant commit, serves 0 stale verdicts whether or not it received the
signal: the recomputed hash differs. An entity-indexed cache that reuses any
verdict for the transaction id behaves identically when notified and returns
the stale verdict in all 8 sequences when not. Flush with a signal is safe but
drops the 8 unrelated verdicts too. The admission process in this system is
unnotified by construction, so the no-signal column is the deployed case, not
a failure injection. These sequences show no property-level selectivity and
no advantage over an entity cache that is always notified.

\begin{figure}[t]
\centering
\includegraphics[width=\columnwidth]{fig_delta_sequence}
\caption{Eight sequences per policy. ``Signal'' is a commit notification. The
dotted line is 8 of 8.}
\Description{Grouped bars for six policy and signal combinations; only the
entity cache without a signal shows stale reuse.}
\label{fig:seq}
\end{figure}

\paragraph{Counters on the live mix}
We keep four events apart: a cold miss (no entry), a dependency-change miss
(an entry with a different hash), retirement (the committing transaction's own
entry removed), and cross-transaction invalidation (another transaction's
entry removed). For 500 offers on one advertisement the ABCI process records
500 cold misses, 500 hits, 500 retirements, and 0 cross-transaction
invalidations: committing one offer does not change the advertisement facts
another offer reads. The reuse is the same offer checked in CheckTx and then
DeliverTx, not sharing across offers.

\emph{Finding.} Both SHACL treatments match the oracle and fail closed.
$\Delta$ served no stale verdict in any tested transition, including without
a commit signal, which is exactly where the entity baseline fails.

\subsection{Validation cost (RQ2)}
\label{sec:cost}

Figure~\ref{fig:phases} times one validation on a validator, outside
consensus. Direct takes 4.3\,ms [4.2, 4.5], a SHACL miss 19.5\,ms [18.9,
20.0], and a $\Delta$ hit 4.6\,ms [4.5, 4.7]. The hit is almost entirely the
4.1\,ms dependency read that recomputes the hash; the lookup itself is
0.004\,ms. A miss spends 4.7\,ms reading dependencies, 0.06\,ms building RDF,
and 13.2\,ms in the HTTP call, of which parsing takes 0.6, SHACL 5.0, and
state checks 3.3\,ms; the rest is transport and JSON. Shapes load once at
startup. Closing the referenced advertisement turns the next lookup into a
miss and a rejection at 22.5\,ms. Phase medians need not sum to the median
total.

\begin{figure}[t]
\centering
\includegraphics[width=\columnwidth]{fig_validation_phases}
\caption{In-process validation, $n{=}200$. Segments are phase medians.
Whiskers are 95\% intervals of the total.}
\Description{Stacked bars for Direct, SHACL miss, Delta hit, and value-change
miss; the hit bar is dependency fetch only.}
\label{fig:phases}
\end{figure}

\paragraph{Ledger size}
Seeding $10^3$, $10^4$, and $10^5$ lifecycles into MongoDB, bypassing
consensus, separates query cost from commit. At $10^5$ the store holds
310{,}000 transaction and metadata documents and 100{,}000 assets. With the
marketplace indexes, Direct stays at 4.9/4.2/4.6\,ms, a miss at
22.5/17.9/19.9\,ms, and a hit at 5.3/4.6/4.5\,ms. Without indexes at $10^5$
the three rise to 540, 1{,}342, and 789\,ms (Figure~\ref{fig:state}). A hit
still reads the projection, so it does not hide a collection scan.

\begin{figure}[t]
\centering
\includegraphics[width=0.9\columnwidth]{fig_state_sweep}
\caption{Seeded ledger, indexes present (solid) and removed (dashed).}
\Description{Log-scale lines; indexed conditions are flat, unindexed
conditions grow with ledger size.}
\label{fig:state}
\end{figure}

\emph{Finding.} A hit removes the sidecar call, about 15\,ms, and costs about
the same as Direct, which checks less. The hit's floor is the dependency read,
and the benefit holds only while dependency queries are indexed.

\subsection{End-to-end behavior (RQ3)}
\label{sec:e2e}

\paragraph{Lifecycle}
A seven-step lifecycle is seven commits. Medians are 4.27\,s [4.20, 4.39]
for Direct, 4.37\,s [4.32, 4.52] for SHACL, and 4.30\,s [4.23, 4.41] for
SHACL+$\Delta$ ($n{=}30$). The intervals overlap and the largest gap is
0.10\,s: low-load parity, not a speedup. An earlier four-region deployment
($n{=}5$) sat near 5.1\,s for every treatment and is not pooled here.

\paragraph{Shared and disjoint offers}
For 500 offers on one advertisement, commit medians are 610\,ms [605, 613],
626\,ms [621, 630], and 572\,ms [569, 575]. $\Delta$ is 54\,ms (8.6\%) below
SHACL in this campaign; the saving is the DeliverTx hit of
Section~\ref{sec:safety}, not sharing across offers. For 200 disjoint offers,
each preceded by its own Create and Advertise, they are 612\,ms [604, 618],
623\,ms [618, 632], and 609\,ms [604, 615], and the intervals overlap.

\paragraph{Offered load}
Eight threads broadcast asynchronously at a fixed rate for 20\,s and do not
wait for commit. Throughput counts commits over wall time including the
drain; latency comes from 16 client samples per rate, because block
timestamps sat at a fixed multi-minute offset. Every submitted transaction
eventually commits. Direct peaks at 37.1\,tx/s when 80\,tx/s is offered.
SHACL and SHACL+$\Delta$ stay between 9.7 and 13.9\,tx/s from 20\,tx/s
onward (Figure~\ref{fig:system}, Table~\ref{tab:open}). At 80\,tx/s the
sampled median latency is 19\,s, 146\,s, and 102\,s. A new transaction is a
cold miss and a hit still reads its dependencies, so $\Delta$ does not raise
the rate for a stream of distinct transactions. These are short bursts, not
steady-state capacity, and each validator runs one sidecar process.

\begin{figure*}[t]
\centering
\includegraphics[width=0.9\textwidth]{fig_system_boundary}
\caption{Left: seven-step lifecycle on the primary deployment ($n{=}30$).
Right: transactions committed per second of wall time, including the drain
after each 20\,s submission window.}
\Description{Left, three bars of equal height with overlapping intervals.
Right, Direct rises to about 37 transactions per second while both SHACL
treatments flatten near 10 to 14.}
\label{fig:system}
\end{figure*}

\begin{table}[t]
\caption{Open-loop sweep. Throughput includes drain. Latency is the median of
16 client samples, in seconds. Below 20\,tx/s offered, all three commit at
the offered rate.}
\label{tab:open}
\small
\setlength{\tabcolsep}{4pt}
\begin{tabular}{@{}r rr rr rr@{}}
\toprule
 & \multicolumn{2}{c}{Direct} & \multicolumn{2}{c}{SHACL} & \multicolumn{2}{c}{SHACL+$\Delta$} \\
\cmidrule(lr){2-3}\cmidrule(lr){4-5}\cmidrule(lr){6-7}
Offered & tx/s & lat. & tx/s & lat. & tx/s & lat. \\
\midrule
20  & 16.0 & 1.4 & 11.1 & 12  & 11.8 & 8.3 \\
40  & 32.0 & 2.3 & 12.9 & 48  & 11.6 & 55 \\
80  & 37.1 & 19  & 10.1 & 146 & 13.9 & 102 \\
160 & 28.9 & 90  & 9.7  & 315 & 10.5 & 290 \\
\bottomrule
\end{tabular}
\end{table}

\emph{Finding.} At low load the three lifecycles are indistinguishable. Under
a stream of distinct transactions, declarative validation cuts achieved
throughput to roughly a third of Direct's at 40 to 80\,tx/s offered, and
$\Delta$ does not remove that first evaluation.

\subsection{Portability to Fabric (RQ4)}
\label{sec:fabric}

On Fabric the chaincode is the authority, and it implements every rule in
the oracle, including the two uniqueness checks Direct lacks on SmartChainDB.
The SHACL treatments use a second, Fabric-adapted shape set in which the
state rules appear as boolean facts the client reads from world state before
endorsement; the client submits only if the pre-check passes, and skipping
the pre-check does not skip the chaincode. All three treatments reject a
non-owner advertisement, Direct inside the chaincode and the SHACL treatments
before endorsement, and all 10 offers placed after close
(Table~\ref{tab:conformance}). Lifecycle medians are 14.31\,s (14.29--14.37),
14.60\,s (14.59--14.77), and 14.57\,s (14.54--14.58), $n{=}5$. The pre-check
median is 31\,ms under SHACL and 26\,ms under SHACL+$\Delta$, about 2\% of
the Direct lifecycle. Fabric does not validate again at endorsement and
commit, so there is no second check to reuse; a repeated driver-side pre-check
is a client memo hit and we do not count it as a speedup. We do not rank the
platforms by throughput. Endorsement, ordering, and commit were not timed
separately, and the seven-type matrix and a sidecar-outage trial were not run
on Fabric.

\emph{Finding.} The rules transfer as re-expressed shapes, and the two
state-dependent rejections agree with the chaincode. The added cost on this
path is the client pre-check, not a cache effect.

\subsection{Limitations}
\label{sec:threats}

Direct omits two uniqueness checks, so its times are not a same-semantics
baseline. The projection is hand-written and hashes whole records. Read-set
coverage was observed, not enforced, when this campaign ran, and the id-level
read set cannot detect a missing absence condition. The in-flight conflict
tracker produced no abort in a racing-close probe and is not claimed. The
live ledger was not reset between treatments. Open-loop windows are 20\,s
with 16 latency samples per rate. All primary numbers come from one campaign
on four validators in one region with one HTTP worker. Fabric is one host,
five lifecycles, and two rejection probes.

\end{document}
